\section*{Chapitre \ref{ch:g3:stones-pebbles-sand} --- Cailloux, galets et sable}

\begin{solution}{exo:g3:stones-pebbles-sand:1}
La \omterm{def:g3:stones-pebbles-sand:rock}{roche en place} est la \omterm{def:g3:stones-pebbles-sand:rock}{roche} solide qui se trouve partout sous le sol.
On la voit là où aucun sol ne la cache~: sur les flancs raides des
\omterm{def:g1:landscapes-around-us:landform}{montagnes}, dans les falaises au bord de la \omterm{def:g2:water-on-earth:ocean}{mer}, le long des routes
taillées à travers une \omterm{def:g1:landscapes-around-us:landform}{colline}.
\end{solution}

\begin{solution}{exo:g3:stones-pebbles-sand:2}
La \omterm{def:g3:stones-pebbles-sand:sand}{boue}, le \omterm{def:g3:stones-pebbles-sand:sand}{sable}, le \omterm{def:g3:stones-pebbles-sand:pebble}{galet}, le bloc.
\end{solution}

\begin{solution}{exo:g3:stones-pebbles-sand:3}
1 millimètre est plus petit que 2 millimètres~: du \omterm{def:g3:stones-pebbles-sand:sand}{sable}. 3 millimètres,
c'est plus que 2 millimètres~: du gravier.
\end{solution}

\begin{solution}{exo:g3:stones-pebbles-sand:4}
Naturels~: le \omterm{def:g3:stones-pebbles-sand:pebble}{galet}, le bloc de granite (taillé par les gens, mais la \omterm{def:g3:stones-pebbles-sand:rock}{roche}
est naturelle), le \omterm{def:g3:stones-pebbles-sand:sand}{sable} de la plage. Fabriqués par les gens~: la brique
et le verre.
\end{solution}

\begin{solution}{exo:g3:stones-pebbles-sand:5}
L'eau de la \omterm{def:g1:landscapes-around-us:water}{rivière} l'a longtemps roulé sur le fond et cogné contre
d'autres cailloux~: ses angles se sont usés.
\end{solution}

\begin{solution}{exo:g3:stones-pebbles-sand:6}
Le granite, pour un mur~; des dalles de pierre minces et plates (une \omterm{def:g3:stones-pebbles-sand:rock}{roche} comme l'ardoise), pour les toits.
\end{solution}

\begin{solution}{exo:g3:stones-pebbles-sand:7}
Le rond~: plus un caillou est roulé longtemps par l'eau, plus ses angles
s'usent. Le pointu n'a pas voyagé loin de la \omterm{def:g3:stones-pebbles-sand:rock}{roche} dont il s'est détaché.
\end{solution}

\begin{solution}{exo:g3:stones-pebbles-sand:8}
Le plus gros \omterm{def:g3:stones-pebbles-sand:pebble}{galet} mesure 3 centimètres. Un centimètre fait
10 millimètres, et $10 \div 2 = 5$~: cinq grains de 2 millimètres tiennent
sur 1 centimètre.
\end{solution}

\begin{solution}{exo:g3:stones-pebbles-sand:9}
Dans la mare tranquille. Une eau rapide garde les minuscules grains de
\omterm{def:g3:stones-pebbles-sand:sand}{boue} en mouvement~; c'est seulement là où l'eau est presque immobile qu'ils
peuvent se déposer au fond.
\end{solution}

\begin{solution}{exo:g3:stones-pebbles-sand:10}
Parce qu'ils sont faits de matériaux tirés du sol~: les briques avec de
l'argile, le béton avec du \omterm{def:g3:stones-pebbles-sand:sand}{sable}, des \omterm{def:g3:stones-pebbles-sand:pebble}{graviers} et du ciment, le verre avec
du \omterm{def:g3:stones-pebbles-sand:sand}{sable}. L'argile, le \omterm{def:g3:stones-pebbles-sand:sand}{sable} et les \omterm{def:g3:stones-pebbles-sand:pebble}{graviers} viennent de \omterm{def:g3:stones-pebbles-sand:rock}{roche} brisée.
\end{solution}
